\documentclass[11pt]{article}

\usepackage[margin=0.9in]{geometry}
\usepackage{amsmath,amssymb}
\usepackage{booktabs}
\usepackage{enumitem}
\usepackage{listings}
\usepackage{xcolor}
\usepackage{hyperref}

\definecolor{backcolour}{rgb}{0.95,0.95,0.92}
\definecolor{codeblue}{rgb}{0.05,0.15,0.55}
\definecolor{codegreen}{rgb}{0.0,0.45,0.0}

\lstset{
    backgroundcolor=\color{backcolour},
    basicstyle=\ttfamily\small,
    keywordstyle=\color{codeblue}\bfseries,
    commentstyle=\color{codegreen},
    stringstyle=\color{purple},
    breaklines=true,
    frame=single,
    showstringspaces=false,
    tabsize=4
}

\title{CPSC 250 \\ Memory Allocation and Garbage Collection}
\author{Edward Brash}
\date{}

\begin{document}

\maketitle

\section{Introduction}

Memory management is one of the most important concepts in computer science.

Many programming languages require the programmer to manually allocate and free memory.

Python attempts to automate much of this process.

However, understanding what is happening internally is still extremely valuable.

\section{Two Philosophies}

There are two common approaches to memory management.

\subsection{Approach 1}

\begin{quote}
Trust Python completely.
\end{quote}

In this approach, we simply allow Python to:

\begin{itemize}
    \item allocate memory,
    \item track objects,
    \item deallocate unused objects.
\end{itemize}

\subsection{Approach 2}

\begin{quote}
Understand the system at a deeper level.
\end{quote}

Even though Python automates memory management, understanding the underlying model helps us:

\begin{itemize}
    \item write better code,
    \item debug difficult problems,
    \item understand performance,
    \item understand why memory leaks occur.
\end{itemize}

\section{The Stack and the Heap}

Programs typically use two major memory regions.

\subsection{The Stack}

The stack is:

\begin{itemize}
    \item highly ordered,
    \item fast,
    \item used for local variables and function calls.
\end{itemize}

The stack behaves somewhat like a stack of plates.

\subsection{The Heap}

The heap is:

\begin{itemize}
    \item less ordered,
    \item more flexible,
    \item used for dynamically allocated objects.
\end{itemize}

The heap behaves more like a pile of objects spread around memory.

\section{CPython and the Virtual Machine}

Python programs are executed by the Python virtual machine.

The most common implementation is:

\[
\text{CPython}
\]

CPython handles much of the memory management automatically.

Languages using similar models include:

\begin{itemize}
    \item Python,
    \item Java,
    \item C\#,
    \item Perl,
    \item Ruby,
    \item Lua.
\end{itemize}

\section{All Python Objects Live on the Heap}

In CPython, essentially all Python objects are stored on the heap.

This includes:

\begin{itemize}
    \item integers,
    \item floats,
    \item strings,
    \item lists,
    \item dictionaries,
    \item user-defined objects.
\end{itemize}

\section{Reference Counting}

CPython keeps track of how many variables refer to each object.

This is called the \textbf{reference count}.

When an object is created:

\begin{itemize}
    \item memory is allocated on the heap,
    \item the object's reference count is initialized.
\end{itemize}

\section{Garbage Collection}

When the reference count of an object falls to zero:

\begin{itemize}
    \item the object is no longer accessible,
    \item its memory can be reclaimed.
\end{itemize}

This process is called garbage collection.

\section{Example Class}

\begin{lstlisting}[language=Python]
class Test:

    def __init__(self, name):
        self.name = name

    def __del__(self):
        print(f"Deleting {self.name}")
\end{lstlisting}

The method:

\begin{lstlisting}[language=Python]
__del__()
\end{lstlisting}

is called when the object is destroyed.

\section{Example Program}

\begin{lstlisting}[language=Python]
if __name__ == "__main__":

    print("Discarded object creation")
    Test("hello1")

    print("Saved object creation")
    my_list = Test("hello2")

    print("Program done")
\end{lstlisting}

\section{What Happens?}

\subsection{Case 1}

\begin{lstlisting}[language=Python]
Test("hello1")
\end{lstlisting}

creates an object on the heap, but no variable stores a reference to it.

Therefore:

\begin{itemize}
    \item reference count quickly falls to zero,
    \item the object becomes garbage collectable almost immediately.
\end{itemize}

\subsection{Case 2}

\begin{lstlisting}[language=Python]
my_list = Test("hello2")
\end{lstlisting}

stores a reference to the object.

Therefore the object remains allocated until the reference disappears.

\section{Reference Counting Example with Strings}

Consider:

\begin{lstlisting}[language=Python]
s1 = "Python"
s2 = "Computer Science"
s3 = s2

s1 = "Zybooks"
s2 = s1
\end{lstlisting}

\section{Step-by-Step Interpretation}

Initially:

\begin{lstlisting}[language=Python]
s1 = "Python"
\end{lstlisting}

creates a string object on the heap with reference count 1.

Then:

\begin{lstlisting}[language=Python]
s2 = "Computer Science"
\end{lstlisting}

creates another string object.

Next:

\begin{lstlisting}[language=Python]
s3 = s2
\end{lstlisting}

causes both \verb|s2| and \verb|s3| to refer to the same object.

The reference count for \verb|"Computer Science"| increases to 2.

\section{Reassignment}

Now consider:

\begin{lstlisting}[language=Python]
s1 = "Zybooks"
\end{lstlisting}

The old string \verb|"Python"| loses its reference.

If no references remain, its reference count becomes zero and it may be garbage collected.

\section{The Garbage Collector Module}

Python includes a garbage collection module:

\begin{lstlisting}[language=Python]
import gc
\end{lstlisting}

We can inspect tracked objects:

\begin{lstlisting}[language=Python]
gc_objects = gc.get_objects()

print(len(gc_objects))
\end{lstlisting}

This can help analyze memory behavior.

\section{Circular References}

Reference counting works well most of the time.

However, circular references can cause problems.

\section{Dog and Cat Example}

Suppose we define:

\begin{lstlisting}[language=Python]
class Dog:

    def __init__(self, name, cousin=None):
        self.name = name
        self.cousin = cousin


class Cat:

    def __init__(self, name, cousin=None):
        self.name = name
        self.cousin = cousin
\end{lstlisting}

Now:

\begin{lstlisting}[language=Python]
bailey = Dog("Bailey")
misty = Cat("Misty", bailey)

bailey.cousin = misty
\end{lstlisting}

creates a circular reference.

\section{The Problem}

Now:

\begin{itemize}
    \item Bailey refers to Misty,
    \item Misty refers to Bailey.
\end{itemize}

Even if we delete one variable:

\begin{lstlisting}[language=Python]
del bailey
\end{lstlisting}

the objects may still refer to each other internally.

Their reference counts may never fall to zero.

\section{Memory Leaks}

This situation can create a memory leak.

A memory leak occurs when memory remains allocated even though the programmer no longer intends to use it.

\section{Breaking the Circular Reference}

One solution is to manually remove the reference:

\begin{lstlisting}[language=Python]
misty.cousin = None
\end{lstlisting}

Now the reference cycle is broken.

\section{Why This Matters}

Programs that repeatedly leak memory may eventually:

\begin{itemize}
    \item slow down,
    \item consume large amounts of RAM,
    \item crash.
\end{itemize}

\section{Key Takeaways}

\begin{itemize}
    \item Python automatically manages most memory allocation.
    \item CPython uses reference counting.
    \item Objects live on the heap.
    \item The stack and heap serve different purposes.
    \item Objects with zero references may be garbage collected.
    \item Circular references can cause memory leaks.
    \item The \verb|gc| module can help inspect memory behavior.
\end{itemize}

\end{document}
